%\documentclass[notes,10pt,aspectratio=169]{beamer}
\documentclass[10pt,aspectratio=169]{beamer}

\usetheme[progressbar=frametitle]{metropolis}
\usecolortheme{rose}
\usecolortheme{default}

\usepackage[utf8]{inputenc}
\usepackage[T1]{fontenc}
\usepackage{lmodern}
\usepackage{xcolor}
\usepackage{tikz}
\usepackage{booktabs}
\usetikzlibrary{shapes.geometric, arrows, positioning}

\tikzstyle{block} = [rectangle, draw, text width=4cm, align=center, rounded corners, minimum height=1cm]
\tikzstyle{decision} = [rectangle, draw, text width=5cm, align=center, fill=blue!10, rounded corners, minimum height=1cm]
\tikzstyle{terminal} = [rectangle, draw, text width=4.5cm, align=center, fill=yellow!30, rounded corners, minimum height=1cm]
\tikzstyle{end} = [rectangle, draw, text width=5cm, align=center, fill=green!30, rounded corners, minimum height=1cm]
\tikzstyle{arrow} = [->, thick]

\usepackage{adjustbox}
\setbeamertemplate{itemize item}[triangle]

\definecolor{accent}{RGB}{78,205,196}
\setbeamercolor{frametitle}{fg=black,bg=white}
\setbeamercolor{item}{fg= orange!80}
\setbeamercolor{button}{bg=orange, fg=white}

%% BibLaTeX disabled (no citations in this presentation)
% \usepackage[backend=biber, style=apa, citestyle=authoryear]{biblatex}
% \addbibresource{../references.bib}

\title{Meeting with Steve}
\author{Lucas Condeza\inst{1}}
\institute{\inst{1} Yale University}
\date{March 26, 2026}

\begin{document}


%-------------------------------
% SLIDE 1: RECALL + TODAY
%-------------------------------
\begin{frame}{Last Meeting (Feb 27)}
\textbf{What we discussed:}
\begin{itemize}
    \item Presented differentiated-products model (logit demand) with switching costs and asymmetric information
    \item Identification strategy: 4-step sequential procedure using $(r_i^{\text{win}}, W_i, D_i)$
    \item To do:  add consumer heterogeneity
\end{itemize}

\vspace{0.3cm}
\textbf{Today's meeting:}
\begin{itemize}
    \item New information structure: bank 1 observes default probability $h$, bank 2 observes noisy signal  $\mathbf{x}$ 
    \item Full identification --- nonparametric, no parametric assumptions on $F(h|\mathbf{x})$ or $G$
\end{itemize}
\end{frame}


%-------------------------------
% SLIDE 2: MODEL SETUP
%-------------------------------
\begin{frame}{Model Setup}

\textbf{Two banks compete for a borrower} 
\begin{center}
\renewcommand{\arraystretch}{1.3}
\begin{tabular}{lcc}
\toprule
& \textbf{Bank~1 (incumbent)} & \textbf{Bank~2 (entrant)} \\
\midrule
Default probability $h$ & Observed & Not observed \\
Covariates $\mathbf{x}\in \{x_1, \ldots, x_K\}$ & \textbf{Not observed} & Observed \\
\bottomrule
\end{tabular}
\end{center}

\begin{itemize}

    \item Logit demand: $s(r_1, r_2) = \frac{1}{1 + \exp(\alpha(r_1 - r_2) - \lambda)}$
\end{itemize}

\vspace{0.2cm}
\textbf{Key consequence:} Bank 1 cannot condition on $\mathbf{x}$, so $\sigma^*(h)$ depends on $h$ alone

$\Rightarrow$ \textbf{One-to-one} mapping from bank-1 rates to borrower types
\end{frame}


%-------------------------------
% SLIDE 3: EQUILIBRIUM
%-------------------------------
\begin{frame}{Equilibrium}

Bank 1 observes $h$, forms posterior $\pi_j(h) = \Pr(\mathbf{x} = x_j \mid h)$ via Bayes' rule, and maximizes expected profit over the $K$ competitive environments:
$$\boxed{\sigma^*(h) = \frac{1}{1-h} + \frac{1}{\alpha} \cdot \frac{\bar{S}(\sigma^*(h);\, h)}{\bar{V}(\sigma^*(h);\, h)}}$$
\begin{itemize}
    \item $\bar{S} = \sum_j \pi_j(h)\, s(\sigma^*(h), r_{2,j}^*)$: expected market share (averaged over competitive environments)
    \item $\bar{V} = \sum_j \pi_j(h)\, s(\sigma^*(h), r_{2,j}^*)(1-s(\sigma^*(h), r_{2,j}^*))$: expected share variance
\end{itemize}

\vspace{0.3cm}
\textbf{Data:} For each borrower $i$, the econometrician observes:
\begin{itemize}
    \item $r_i^{\text{win}}$: winning interest rate
    \item $W_i \in \{1, 2\}$: identity of the winning bank
    \item $D_i \in \{0, 1\}$: default indicator
\end{itemize}
\end{frame}


%-------------------------------
% SLIDE 4: TYPE RECOVERY + NONPARAMETRIC OBJECTS
%-------------------------------
\begin{frame}{Identification: Type Recovery and Nonparametric Objects}

Since $\sigma^*(h)$ is one-to-one, among bank-1 winners:
$$E[D_i \mid r_i^{\text{win}} = r_1,\; W_i = 1] = \tau(r_1) = h$$

\begin{center}
\renewcommand{\arraystretch}{1.3}
\begin{tabular}{ll}
\toprule
\textbf{Object} & \textbf{Identified from} \\
\midrule
$\sigma^*(h)$, $\tau(r_1)$ & $E[D \mid r_1, W\!=\!1]$ (type recovery) \\
$m(h) = \sigma^*(h) - \frac{1}{1-h}$ & Directly from $\sigma^*(h)$ \\
$r_{2,j}^*$ (distinct bank-2 rates) & Bank-2 winning rates take $K$ distinct values \\
$E[D \mid r_2, W\!=\!2]$ & Regression among bank-2 winners \\
\bottomrule
\end{tabular}
\end{center}

\vspace{0.2cm}
\textbf{What remains:} $(\alpha, \lambda)$, population probabilities $(p_j)$, type distributions $f(h|x_j)$
\end{frame}


%-------------------------------
% SLIDE 5: TWO POINTWISE RESTRICTIONS
%-------------------------------
\begin{frame}{Identification: Two Restrictions at Each $h$}

Define $\gamma_j(h) \equiv p_j\, f(h|x_j)$: the joint density of $(x_j, h)$  and $s_j(h) \equiv s(\sigma^*(h), r_{2,j}^*;\, \alpha, \lambda)$.  At each~$h$, two restrictions:

\vspace{0.15cm}
\textbf{(R1) First-order condition:}
$$\sum_{j=1}^K \gamma_j(h)\, s_j(h)\Big[1 - \alpha\, m(h)\,\big(1\!-\!s_j(h)\big)\Big] = 0$$

\textbf{(R2) Law of total probability for bank-1 winners:}
$$\sum_{j=1}^K \gamma_j(h)\, s_j(h) = f(h,\, W_i\!=\!1)$$
Right-hand side is data: type recovery gives $h_i = \tau(r_i^{\text{win}})$ for each bank-1 winner, so the empirical distribution of~$h$ among stayers directly estimates the joint density $f(h,\, W_i\!=\!1)$.

\vspace{0.15cm}
$\Rightarrow$ \textbf{2 equations in $K$ unknowns} $\gamma_1(h), \ldots, \gamma_K(h)$ at each $h$
\begin{itemize}
    \item $K = 2$: exactly identified pointwise $\Rightarrow$ \textbf{nonparametric} identification
    \item $K \geq 3$: underdetermined pointwise $\Rightarrow$ parametric family needed
\end{itemize}
\end{frame}


%-------------------------------
% SLIDE 6: K>=3 CASE
%-------------------------------
\begin{frame}{Case $K \geq 3$: Parametric Identification}

2 equations, $K$ unknowns per $h$ $\Rightarrow$ $(K-2)$ free functions $\Rightarrow$ nonparametric ID fails.


\vspace{0.2cm}
\textbf{Resolution:} parametrize $f(h|x_j) = f(h;\, \theta_j)$ with a $P$-parameter family (e.g., Beta: $P = 2$).

\vspace{0.2cm}
(R2) becomes: $\sum_{j=1}^K p_j\, f(h;\theta_j)\, s_j(h) = f(h,\, W_i\!=\!1)$ for all $h$

\vspace{0.2cm}
\textbf{Identification (same sequential logic as $K = 2$):}
\begin{enumerate}
    \item Given $(\alpha, \lambda)$: continuum of equations pins down $\underbrace{KP}_{\substack{P \text{ params per} \\ \text{density} \times K \text{ groups}}} + \underbrace{K - 1}_{\substack{p_1,\ldots,p_{K-1} \\ (p_K = 1 - \sum p_j)}}$ scalar unknowns
    \item $(\alpha, \lambda)$ pinned down by the $K$ bank-2 default-rate conditions ($K \geq 3 > 2$ unknowns)
\end{enumerate}

\end{frame}


%-------------------------------
% SLIDE 8: REDUCED-FORM EVIDENCE (REDEC)
%-------------------------------
\begin{frame}{Reduced-Form Evidence: REDEC Policy (Chile, April 2026)}

\textbf{Setting:} Chile creates a centralized debt registry (REDEC) --- all lenders must share positive \& negative credit info.  Non-bank borrowers' histories become visible to banks for the first time.

\vspace{0.2cm}
\textbf{Regressions} (sample: non-bank-only borrowers, $Good_i = \mathbf{1}[\text{no delinquencies}]$):

\vspace{0.1cm}
\textbf{1.\ Credit access:}
$$NewBank_{it} = \beta^A \cdot (Post_t \times Good_i) + \gamma X_{it} + \delta_i + \theta_t + \varepsilon_{it}$$
Expected: $\hat{\beta}^A > 0$ --- good types gain bank access once histories are visible

\textbf{2.\ Pricing:}
$$r_{ijt} = \beta^P \cdot (Post_t \times Good_i) + \gamma X_{it} + \mu_j + \theta_t + \varepsilon_{ijt}$$
Expected: $\hat{\beta}^P < 0$ --- good types get lower rates (separating vs.\ pooling)

\textbf{3.\ Delinquency:}
$$Delinq_{ijt} = \beta^D \cdot (Post_t \times OCNB\_Debt_i) + \gamma X_{it} + \mu_j + \theta_t + \varepsilon_{ijt}$$
Expected: $\hat{\beta}^D < 0$ --- banks screen on non-bank debt, improving loan quality
\end{frame}


%==================================
% APPENDIX
%==================================
 
\begin{frame}{Appendix: Case $K = 2$ --- Nonparametric Identification}

\textbf{Step 1:} Given candidate $(\alpha, \lambda)$, solve (R1)--(R2) --- a $2 \times 2$ system at each~$h$.

Since $r_{2,1}^* \neq r_{2,2}^*$ $\Rightarrow$ $s_1(h) \neq s_2(h)$ $\Rightarrow$ system is nonsingular
$\Rightarrow$ $\gamma_j(h;\, \alpha, \lambda)$ determined at every $h$.

\vspace{0.2cm}
\textbf{Step 2:} Recover $p_j$: integrability $\int f(h|x_j)\, dh = 1$ gives $p_j = \int \gamma_j(h)\, dh$.

\vspace{0.2cm}
\textbf{Step 3:} Pin down $(\alpha, \lambda)$ using bank-2 conditional default rates:
$$E[D_i \mid r_{2,j}^*,\, W_i = 2] = \frac{\int h\,(1 - s_j)\, f(h|x_j)\, dh}{\int (1 - s_j)\, f(h|x_j)\, dh}, \quad j = 1, 2$$
Two equations, two unknowns $\Rightarrow$ $(\alpha, \lambda)$ identified.
\end{frame}


\end{document}
